\chapter{Introducción}\label{cap:1}

\section{Motivación y contexto}\label{sec:motivacion}

Chile central ha atravesado desde 2010 una megasequía sin precedentes en el registro instrumental, acompañada de una marcada variabilidad hidrológica interanual (Garreaud et al., 2017, 2020; Boisier et al., 2016, 2018). Las métricas agregadas de precipitación no capturan por completo cómo se organizan en el tiempo los períodos secos: dos climas con el mismo acumulado anual ---por ejemplo, 300 mm en dos eventos frente a quince--- pueden tener consecuencias muy distintas para embalses, agricultura y ecosistemas, según cuán prolongadas sean las rachas secas. La pregunta no es exclusiva de Chile central: el país abarca desde el desierto de Atacama, con rachas secas de meses o años, hasta el sur húmedo, donde las rachas son mucho más cortas, lo que permite estudiar la persistencia seca a lo largo de un gradiente climático extremo.

Este trabajo estudia esa persistencia mediante dry spells: rachas consecutivas de días secos, definidos como días con precipitación por debajo de un umbral de día húmedo. En lugar de preguntar cuánto llueve en promedio, el enfoque distributivo pregunta con qué frecuencia ocurren las rachas y qué tan largas llegan a ser, incluida la cola de eventos extremadamente prolongados. Las PDFs de duración ofrecen un marco parsimonioso para caracterizar esa intermitencia, permitiendo distinguir cambios en la frecuencia total de rachas, en la forma de la distribución y, particularmente, en su cola (Stechmann y Neelin, 2014; Martinez-Villalobos y Neelin, 2018, 2019).

En la literatura climática, la persistencia seca se ha cuantificado sobre todo con índices escalares. El más difundido es el de días secos consecutivos (CDD), definido por el ETCCDI como el máximo anual o estacional de días seguidos con precipitación inferior a 1 mm (Zhang et al., 2011; Sillmann et al., 2013a, 2013b). Otros trabajos reportan el número de días secos o la duración media de las rachas (Giorgi et al., 2011, 2019; Polade et al., 2014). Un conjunto más reducido examina la distribución completa de duraciones: Zolina et al.~(2013) mostraron que los cambios en rachas húmedas y secas en Europa no se reducen a un cambio en el número de días húmedos, sino a una reorganización de su agrupación temporal; Lana et al.~(2006) ajustaron distribuciones a dry spells largos en la Península Ibérica. En Chile, ajustes paramétricos confirman colas pesadas y desviaciones de un proceso sin memoria (Herrera-Marín et al., 2026). Estas distinciones importan porque un mismo aumento de días secos puede fragmentarse en muchas rachas cortas o concentrarse en pocas rachas muy largas. El CDD retiene el extremo de la cola en cada año, pero no describe con qué frecuencia ocurren spells de 10, 20 o 50 días ni cómo se desplaza esa cola (Trenberth et al., 2003; Pendergrass y Knutti, 2018).

A ese panorama se suman las proyecciones de menor precipitación en Chile central bajo calentamiento, con señal robusta en CMIP6 hacia fines de siglo, incluidas trayectorias de altas emisiones como SSP5-8.5 (Araya-Osses et al., 2020; Salazar et al., 2024). A escala global, CMIP5/CMIP6 proyectan más días secos y dry spells más largos en subtrópicos y climas mediterráneos, Chile incluido: el aumento de días secos explica gran parte del secado proyectado (Polade et al., 2014, 2017; Sillmann et al., 2013b; Giorgi et al., 2019; Wainwright et al., 2021). Saber que ``llueve menos en promedio'' no equivale, sin embargo, a saber cómo cambia la distribución de duraciones de las rachas secas. Un mismo secado puede expresarse como muchas rachas relativamente cortas o como menos interrupciones y rachas mucho más persistentes, y solo la distribución completa, en particular su cola, permite distinguir entre ambas formas. La pregunta central es, por tanto, cómo cambia la distribución de dry spells bajo calentamiento y, en especial, la frecuencia de rachas muy largas.

En Chile, el diagnóstico del secado reciente se ha construido sobre anomalías y tendencias de precipitación y caudales (Boisier et al., 2016, 2018; Garreaud et al., 2017, 2020), y las proyecciones, sobre acumulados o impactos hidrológicos (Araya-Osses et al., 2020; Bozkurt et al., 2018; Salazar et al., 2024). Índices estandarizados (SPI, SPEI, SSI) caracterizan el déficit en escalas temporales, pero no la reorganización día a día de las rachas (Báez-Villanueva et al., 2024). Herrera-Marín (2026) y Herrera-Marín et al.~(2026) modelan, con CR2MET, los inicios de dry spells y su agrupación, pero no el cambio de la distribución de duraciones bajo calentamiento. El antecedente más cercano es Latoja et al.~(2025), quienes, con CR2MET, encuentran que entre 1960 y 2021 los dry spells en Chile central (32°--40°S) se volvieron más numerosos y más cortos en promedio, con cambios significativos sobre todo al sur de 36°S: la precipitación se volvió más intermitente. Esa tendencia contrasta con las rachas más largas que los modelos globales proyectan para climas mediterráneos, más días secos y rachas más largas. Queda abierto, por tanto, si el secado futuro de Chile repetirá la mayor intermitencia documentada en el período observado o si se reorganizará hacia menos interrupciones y una mayor frecuencia de rachas muy largas.

Responder esa pregunta con un modelo regional exige, primero, saber si el modelo representa la persistencia histórica. Los modelos tienden a llover demasiado a menudo y con poca intensidad (Stephens et al., 2010; Martinez-Villalobos et al., 2022), lo que fragmenta las rachas secas. La mayoría de los modelos CMIP5 y CMIP6 subestima la racha seca anual más larga observada, y los que más la subestiman tienden a proyectar aumentos futuros menores (Petrova et al., 2024). Comparar persistencia entre observaciones y modelos exige, por tanto, controlar la frecuencia de días húmedos antes de interpretar diferencias de duración (Leung et al., 2022), y evaluar la cola histórica es condición para juzgar la credibilidad y la magnitud del cambio proyectado.

En este estudio, centrado en Chile continental y con especial relevancia para Chile central, examinamos cómo cambia la distribución de duraciones de dry spells hacia mediados de siglo bajo SSP5-8.5 y, en particular, la frecuencia de rachas prolongadas. Como paso previo para interpretar esa señal futura, evaluamos en qué medida el modelo regional ALADIN reproduce la distribución histórica respecto del producto grillado CR2MET, una vez controladas las diferencias en la frecuencia de días húmedos. Así abordamos no solo cuánto cambia la precipitación, sino cómo se reorganiza temporalmente la persistencia de los períodos secos.

\section{Preguntas de investigación, objetivos e hipótesis (borrador)}\label{sec:objetivos}

\textbf{Preguntas de investigación}

\begin{enumerate}
\def\labelenumi{\arabic{enumi}.}
\tightlist
\item
  Una vez igualada la frecuencia de días húmedos, ¿reproduce el modelo regional ALADIN la distribución histórica de duraciones de dry spells observada en CR2MET a lo largo de Chile continental?
\item
  Hacia mediados de siglo (2040--2074), ¿cómo cambian la distribución de duraciones y, en particular, la frecuencia de rachas secas largas?
\end{enumerate}

\textbf{Objetivo general}

Caracterizar la distribución de duraciones de dry spells en Chile continental en el clima actual y su cambio proyectado a mediados de siglo, evaluando previamente la capacidad del modelo regional ALADIN para reproducirla.

\textbf{Objetivos específicos}

\begin{enumerate}
\def\labelenumi{\arabic{enumi}.}
\tightlist
\item
  Caracterizar la climatología observada de duraciones de dry spells (CR2MET, 1980--2014) a lo largo del gradiente climático norte--sur.
\item
  Evaluar la representación de esas duraciones en ALADIN frente a CR2MET, tras calibrar el umbral de día húmedo, y cuantificar el sesgo en la duración media y en la cola.
\item
  Cuantificar el cambio proyectado (2040--2074, SSP5-8.5) en la frecuencia de rachas largas mediante risk ratios y en la forma de las distribuciones de duración.
\item
  Evaluar la robustez de los resultados frente a la definición del umbral, el método de regrillado y el escenario de emisiones.
\end{enumerate}

\textbf{Hipótesis (propuesta)}

H1. Aun con la frecuencia de días húmedos igualada, ALADIN subestima la persistencia seca observada, especialmente en la cola de la distribución.

H2. Hacia mediados de siglo aumenta el número de rachas secas largas, aunque el número total de rachas disminuya.

\pendiente{confirmar si el programa exige hipótesis formales; si es así,
justificarlas con la literatura del Capítulo 2.}

\section{Estructura de la tesis}\label{sec:estructura}

El Capítulo 2 presenta el marco teórico y el área de estudio. El Capítulo 3 describe los datos y la metodología. El Capítulo 4 evalúa la representación histórica de las rachas secas en ALADIN, y el Capítulo 5 analiza su cambio proyectado. El Capítulo 6 discute los resultados y sus limitaciones, y el Capítulo 7 presenta las conclusiones. Los anexos reúnen las pruebas de sensibilidad y la información de código y datos.
